% QCM
\begin{enumerate}
    \item Le travail de la force $\vec{F}$, de valeur $F=\qty{5.0}{\N}$, pour
          le déplacement $AB$ (de $A$ à $B$) tel que $AB=\qty{50}{\cm}$
          vaut~:

          \begin{tcolorbox}
          \begin{enumerate*}[itemjoin={\hfill}]
              \item
                    \begin{tikzpicture}[
                            baseline={([yshift=-3mm]current bounding box.north west)},
                        point/.style={circle,inner sep=1.4pt,fill}]
                      \foreach \p/\y in {A/0,B/-1.5}
                      \node[point,label=left:$\p$] (\p) at (0,\y) {};
                      \draw[thick,postaction={decorate},
                            decoration={markings,mark=at position 0.57 with
                            {\arrow{Latex}}}]
                        (A.center) -- (B.center);
                      \draw[->,red,very thick]
                        (A.center) -- node[right] {$\vec{F}$} ($(A)!0.4!(B)$);
                    \end{tikzpicture}
              \item
                    \begin{tikzpicture}[
                        baseline={([yshift=-3mm]current bounding box.north west)},
                        point/.style={circle,inner sep=1.4pt,fill}]
                      \foreach \p/\a in {A/0,B/1.5}
                      \node[point,label={[below left,inner sep=2pt]:$\p$}]
                        (\p) at (-45:\a) {};
                      \draw[thick,postaction={decorate},
                            decoration={markings,mark=at position 0.57 with
                            {\arrow{Latex}}}]
                        (A.center) -- (B.center);
                      \draw[->,red,very thick]
                        (A.center) -- node[above right,inner sep=1pt]
                        {$\vec{F}$} ($(A)!-0.4!(B)$);
                    \end{tikzpicture}
              \item
                    \begin{tikzpicture}[
                        baseline={([yshift=-3mm]current bounding box.north west)},
                        point/.style={circle,inner sep=1.4pt,fill}]
                      \foreach \p/\x in {A/0,B/1.5}
                      \node[point,label=below:$\p$] (\p) at (\x,0) {};
                      \draw[thick,postaction={decorate},
                            decoration={markings,mark=at position 0.57 with
                            {\arrow{Latex}}}]
                        (A.center) -- (B.center);
                      \draw[->,red,very thick]
                        (A.center) -- node[left] {$\vec{F}$} ++(0,1);
                      \draw[thick] (0,2mm) -| (2mm,0);
              \end{tikzpicture}
          \end{enumerate*}
          \end{tcolorbox}
          \begin{tabular}{llll} 
            Cas a~:  & $\square\qty{250}{\J}$~; & $\square\qty{2.5}{\J}$~;
                    & $\square\qty{10}{\J}$. \\
            Cas b~:  & $\square\qty{2.5}{\J}$~; & $\square\qty{250}{\J}$~;
                    & $\square\qty{-2.5}{\J}$. \\
            Cas c~:  & $\square\qty{2.5}{\J}$~; & $\square\qty{0}{\J}$~;
                    & $\square\qty{-2.5}{\J}$. \\
          \end{tabular}
    \item La puissance du travail de la force effectuant un travail de
          $\qty{3000}{\J}$ en 1,0 minute est~: \\
          \begin{itemize*}[label=$\square$,itemjoin=\quad]
              \item $\qty{3000}{\W}$~;
              \item $\qty{50}{\W}$~;
              \item $\qty{50}{\kW}$.
          \end{itemize*}
\end{enumerate}

